\subsection{INTERIOR, CLOSURE, FRONTIER OF A SET; DENSE SETS}

DEFINITION 9. In a topological space X, a point x is said to be an interior point of a subset A of X if A is a neighbourhood of x. The set of interior points of A is called the interior of A and is denoted by \(\mathring{A}\).

According to Definitions 9 and 4, a point x is an interior point of A if there is an open set contained in A which contains x; it follows that \(\mathring{A}\) is the union of all the open sets contained in A, and hence is the largest open set contained in A: in other words, if B is an open set contained in A, then B ⊂ \(\mathring{A}\). Consequently, if A and B are two subsets of X such that B ⊂ A, then \(\mathring{B} \subset \mathring{A}\); and A is a neighbourhood of B if and only if B ⊂ \(\mathring{A}\).

Remark. The interior of a non-empty set can be empty; this is the case for a set consisting of a single point which is not open, for example on the rational line * (or the real line) *.

Proposition 1 of no. 2 can be restated as follows:

A set is open if and only if it coincides with its interior.

The property \((V_{II})\) of no. 2 implies that every point which is an interior point of each of two subsets A and B is an interior point of \(A \cap B\) ; consequently

\[
\mathring {\mathbf {A} \cap \mathbf {B}} = \mathring {\mathbf {A}} \cap \mathring {\mathbf {B}}.\tag{1}
\]

Every point which is interior to the complement of a set A is said to be an exterior point of A, and the set of these points is called the exterior of A in X; a point \(x \in X\) which is an exterior point of A is therefore characterized by the property that x has a neighbourhood which does not meet A.

DEFINITION 10. The closure of a subset A of a topological space X is the set of all points \(x \in X\) such that every neighbourhood of x meets A, and is denoted by \(\overline{A}\) .

This definition can be reformulated by saying that a point x lies in the closure of a set A if there are points of A as near x as we please to.

Every point which is not in the closure of A is exterior to A, and conversely; thus we have the formulae (which are duals of each other)

\[
\mathbf {C} \overline {{\mathbf {A}}} = \mathring {\mathbf {C} \mathbf {A}}, \quad \mathbf {C} \mathring {\mathbf {A}} = \overline {{\mathbf {C} \mathbf {A}}}.\tag{2}
\]

Hence, to any proposition on interiors of sets, there corresponds by duality a proposition on closures, and vice versa. In particular, the closure of a set A is the smallest closed set which contains A; in other words, if B is a closed set such that \(A \subset B\) , then \(\overline{A} \subset B\) . If A and B are two subsets of X such that \(A \subset B\) , then \(\overline{A} \subset \overline{B}\) .

A set is closed if and only if it coincides with its closure.

The dual of formula (1) is

\[
\overline {{{\mathbf {A} \cup \mathbf {B}}}} = \overline {{{\mathbf {A}}}} \cup \overline {{{\mathbf {B}}}}.\tag{3}
\]

PROPOSITION 5. Let A be an open set in X; then for every subset B of X we have

\[
\mathbf {A} \cap \overline {{\mathbf {B}}} \subset \overline {{\mathbf {A} \cap \mathbf {B}}}.\tag{4}
\]

For suppose \(x \in A \cap \overline{B}\) ; then if V is any neighbourhood of x, \(V \cap A\) is a neighbourhood of x, since A is open; hence \(V \cap A \cap B\) is not empty and therefore x lies in the closure of \(A \cap B\) .

If \(x\) lies in the closure of \(A\) but not in \(A\) , then every neighbourhood of \(x\) contains a point of \(A\) other than \(x\) ; but if \(x \in A\) it can happen that \(x\) has a neighbourhood which contains no point of \(A\) except \(x\) . We say then that \(x\) is an isolated point of \(A\) . In particular, \(x\) is isolated in the whole space \(X\) if and only if \(\{x\}\) is an open set.

A closed set which has no isolated points is called a perfect set.

DEFINITION 11. In a topological space X, a point x is said to be a frontier point of a set A if x lies in the closure of A and in the closure of CA. The set of frontier points of A is called the frontier of A.

The frontier of A is therefore the set \(\overline{A} \cap \overline{\complement A}\) , which is closed. A frontier point x of A is characterized by the property that every neighbourhood of x contains at least one point of A and at least one point of CA; x may or may not belong to A. The frontier of A is the same as the frontier of CA. The interior of A, the exterior of A and the frontier of A are mutually disjoint and their union is the whole space X.

DEFINITION 12. A subset A of a topological space X is said to be dense in X (or simply dense, if there is no ambiguity about X) if \(\overline{A} = X\) , i.e.~if every non-empty open set U of X meets A.

Examples. * We shall see in Chapter IV, § 1 that the set of rational numbers and its complement are dense on the real line. *

In a discrete space X the only dense subset of X is \(X^{*}\) itself. On the other hand, every non-empty subset of X is dense in the topology on X for which the only open sets are \(\varnothing\) and X.

PROPOSITION 6. If \(\mathfrak{B}\) is a base of the topology of a topological space \(X\) , there is a dense set \(D\) in \(X\) such that \(\operatorname{Card}(D) \leqslant \operatorname{Card}(\mathfrak{B})\) .

We may restrict ourselves to the case in which none of the sets of \(\mathfrak{B}\) is empty (the non-empty sets of \(\mathfrak{B}\) already form a base of the topology of \(\mathbf{X}\) ). For each \(\mathbf{U} \in \mathfrak{B}\) , let \(x_{\mathbf{U}}\) be a point of \(\mathbf{U}\) ; it follows from Proposition 3 of no. 3 that the set \(\mathbf{D}\) of the points \(x_{\mathbf{U}}\) is dense in \(\mathbf{X}\) , and we have Card \((\mathbf{D}) \leqslant\) Card \((\mathfrak{B})\) (Set Theory, Chapter III, § 3, no. 2, Proposition 3).

